%% AFRITIÐ ÞENNAN TEXTA Í YKKAR EIGIÐ SKJAL Í OVERLEAF %%
%% Leiðbeiningar um hvernig á að gera hitt og þetta má finna í skjalinu:
%% https://www.overleaf.com/read/whctggtbjghw
%% Höfundur: Guðný Halldórsdóttir

\documentclass[a4paper]{article}
\usepackage[icelandic]{babel}
\usepackage[utf8x]{inputenc}
\usepackage[T1]{fontenc}
\usepackage[a4paper,top=3cm,bottom=2cm,left=3cm,right=3cm,marginparwidth=1.75cm]{geometry} 
\usepackage{float}
\usepackage{graphicx}
\usepackage{caption}
\usepackage{wrapfig}
% til að setja hlekki (linka)
\usepackage[colorlinks=False]{hyperref}

\usepackage{amsmath, amsfonts, amssymb}  % stærðfræðitákn
\usepackage{color} % til að geta litað texta
\usepackage{enumitem} % fyrir lista og upptalningar
\usepackage{minted} % til að geta sett in kóða

\author{Nafn þitt} %
\title{STÆ205G: Stærðfræðigreining II \\ Vor 2020: Skiladæmi n}
\date{\today}
\begin{document}
\maketitle
\section*{Dæmi 1}
Hér skrifum við fyrirmæli verkefnisins

\subsection*{Lausn:}
Hér kemur þín lausn að verkefninu.

\section*{Dæmi 2}
Hvað er $14\cdot 29$?

\subsection*{Lausn}
Hægt er að reikna dæmið með því að slá því í reiknivél eða:
$$14 \cdot 29 = (10+4)\cdot(30-1)=10\cdot 30 +4\cdot 30 +10\cdot(-1)+4\cdot(-1) = 300+120-10-4= \underline{406} $$
\end{document}
